%! TeX root: ../article.tex
% Anexo de materiales: figuras descriptivas del conjunto de datos y tablas de
% catálogo de componentes. Trasladado del cuerpo al recortar a ~20pp.

\section{Conjunto de datos y catálogo de componentes}
\label{anexo:materiales}

\subsection{Figuras complementarias del conjunto de datos}
\label{anexo:datos-figuras}

\begin{figure}[h!]
  \centering
  \includegraphics[width=\linewidth]{figures/dataset-evolution.pdf}
  \caption[Progreso de la toma de imágenes de grafiti]{Progreso de toma de imágenes
  de grafiti a lo largo del tiempo. La gran mayoría de las imágenes se
capturaron en el otoño de 2022. Durante el primer trimestre de 2023 se
añadieron las últimas 500 imágenes al conjunto de datos.}
  \label{fig:dataset-capture-evolution}
\end{figure}

\begin{figure}[h!]
  \centering
  \includegraphics[width=\linewidth]{figures/dataset_style_distribution.pdf}
  \caption[Distribución por clase de los recortes etiquetados]{Distribución por clase de los recortes etiquetados por estilo,
  separando los subconjuntos de entrenamiento y evaluación.}
  \label{fig:dataset-style-distribution}
\end{figure}

\subsection{Catálogo de componentes}
\label{anexo:catalogo}
\label{subsec:configuraciones}
Esta sección cataloga los componentes concretos integrados en el sistema, que
constituyen los materiales sobre los que opera la evaluación. Su combinación
parametrizada para definir cada experimento se describe en
\ref{subsubsec:espacio-configs}.

El sistema integra siete modelos de extracción de características, todos
empleados congelados: las CNN preentrenadas sobre ImageNet ResNet50, VGG16,
InceptionV3 y MobileNetV3, los \textit{transformers}
DINOv2~\cite{oquab_dinov2_2024} y CLIP~\cite{radford_learning_2021}, y la red de
detección YOLO~\cite{ultralytics_ultralytics_2024}. En el agrupamiento se
cubren los principales paradigmas (centroides, densidad, jerárquicos,
probabilísticos y espectrales) con ocho algoritmos, y en la reducción de
dimensionalidad cinco técnicas (PCA, Kernel PCA, Isomap, UMAP y la identidad,
que no transforma y permite medir el efecto de añadir la etapa). Las
tablas~\ref{tab:embedding-models}, \ref{tab:clustering-algos}
y~\ref{tab:reduction} detallan cada opción con su arquitectura, dimensionalidad,
hiperparámetros y referencias. Las implementaciones de agrupamiento y reducción
provienen de scikit-learn~\cite{pedregosa_scikit-learn_2011}; algunos algoritmos
(K-medias, DBSCAN, GMM) estiman su hiperparámetro principal cuando no se
declara, según el procedimiento de \ref{subsubsec:hiperparam-auto} (columna
«Auto. $k$» de la tabla~\ref{tab:clustering-algos}).

\begin{longtable}{|p{3.2cm}|p{3.8cm}|c|p{4.2cm}|} \hline
  \textbf{Modelo} & \textbf{Arquitectura} & \textbf{Dim.} & \textbf{Preentrenamiento} \\ \hline
  ResNet50~\cite{he_deep_2015} & CNN residual & $2048$ & ImageNet \\ \hline
  VGG16~\cite{simonyan_very_2015} & CNN profunda & $4096$ & ImageNet \\ \hline
  InceptionV3~\cite{szegedy_rethinking_2015} & CNN multi-rama & $2048$ & ImageNet \\ \hline
  MobileNetV3~\cite{howard_searching_2019} & CNN ligera & $1280$ & ImageNet \\ \hline
  DINOv2 ViT-S/14~\cite{oquab_dinov2_2024} & \textit{Transformer}~\cite{dosovitskiy_image_2021} & $384$ & Autosupervisado (LVD-142M) \\ \hline
  CLIP ViT-B/32~\cite{radford_learning_2021} & \textit{Transformer} & $512$ & Imagen-texto (WIT) \\ \hline
  YOLO11 (n/m)~\cite{ultralytics_ultralytics_2024} & Detector \textit{one-stage} & $256$--$512$ & COCO \\ \hline
  \caption{Modelos de extracción de características integrados.}
  \label{tab:embedding-models}
\end{longtable}

\begin{longtable}{|l|l|c|} \hline
  \textbf{Algoritmo} & \textbf{Hiperparámetros principales} & \textbf{Auto. $k$} \\ \hline
  K-medias~\cite{arthur_k-means_2007} & $n\_clusters$ & sí \\ \hline
  DBSCAN~\cite{ester_density-based_1996} & $\epsilon$, $min\_samples$ & sí ($\epsilon$) \\ \hline
  HDBSCAN~\cite{campello_density-based_2013} & $min\_cluster\_size$, $max\_cluster\_size$ & no requiere \\ \hline
  OPTICS~\cite{ankerst_optics_1999} & $min\_samples$, $max\_eps$, $metric$ & no requiere \\ \hline
  Aglomerativo~\cite{mullner_modern_2011} & $n\_clusters$, $linkage$ & no \\ \hline
  Espectral~\cite{ng_spectral_2001,huang_ultra-scalable_2020} & $n\_clusters$, afinidad & no \\ \hline
  GMM~\cite{dempster_maximum_1977} & $n\_components$, covarianza & sí \\ \hline
  \textit{Affinity Propagation}~\cite{frey_clustering_2007} & $damping$, $preference$ & no requiere \\ \hline
  \caption{Algoritmos de agrupamiento implementados.}
  \label{tab:clustering-algos}
\end{longtable}

\begin{longtable}{|l|l|l|} \hline
  \textbf{Técnica} & \textbf{Tipo} & \textbf{Parámetros principales} \\ \hline
  Identity & --- & --- \\ \hline
  PCA~\cite{jolliffe_principal_2002} & Lineal & $n\_components$ \\ \hline
  Kernel PCA~\cite{scholkopf_nonlinear_1998} & No lineal & $n\_components$, \textit{kernel} \\ \hline
  Isomap~\cite{tenenbaum_global_2000} & No lineal & $n\_components$, $n\_neighbors$ \\ \hline
  UMAP~\cite{mcinnes_umap_2020} & No lineal & $n\_components$, $n\_neighbors$, $min\_dist$ \\ \hline
  \caption{Técnicas de reducción de dimensionalidad implementadas.}
  \label{tab:reduction}
\end{longtable}

\subsubsection{Sistemas de almacenamiento}
La persistencia de los \textit{embeddings} y los metadatos asociados se abstrae
tras una interfaz común, de la que se han implementado dos \textit{backends}
con características operativas distintas. El primero, basado en SQLite, no
requiere servicios externos, almacena los vectores como \textit{blobs} y
realiza la búsqueda por similitud mediante un recorrido lineal; resulta
apropiado para despliegues ligeros y como referencia de cálculo exacto frente a
alternativas aproximadas. El segundo se apoya en la extensión \textit{pgvector}
de PostgreSQL y delega la búsqueda al motor de la base de datos. Por defecto la
consulta es un recorrido exacto; opcionalmente puede habilitarse un índice
aproximado HNSW~\cite{malkov_efficient_2018} sobre los vectores, lo que permite
contrastar el coste de la búsqueda exacta frente a la aproximada. El índice
HNSW de \textit{pgvector} impone un límite de 2000 dimensiones sobre el tipo
\texttt{vector}, por lo que no es aplicable a los modelos cuya salida lo supera
(ResNet50 e InceptionV3, con 2048, y VGG16, con 4096); para estos casos debe
recurrirse a una reducción de dimensionalidad previa al almacenamiento o al
\textit{backend} SQLite con búsqueda exacta. Ambos \textit{backends} comparten
el esquema lógico y la persistencia de metadatos, lo que permite cambiar de uno
a otro sin modificar el resto del \textit{pipeline}.
